% introduction.tex
% Foundations-toolkit introduction.  Replaces the template's exemplar
% introduction (which demonstrated cross-reference / figure-include
% patterns).  Reference patterns kept; exemplar prose removed.

\label{sec:introduction}

\placeholder{Rewrite once Phase 2 sections land.  This stub fills the
role of an introduction at v0.1: states the problem, the contribution,
and the scope, and surfaces the audit table.}

\subsection{Problem}
\label{subsec:problem}

The A=1 Discrete Causal Lattice (DCL) programme establishes a
bipartite octahedral lattice in (3+1) dimensions as the substrate of
fundamental physics~\cite{example2024}.\footnote{Citations through
this draft refer to placeholder BibTeX entries in
\texttt{paper-bib/references.bib}; real cites land as the prose is
written.}  Paper~I sets out the substrate-priority claim in 3-D
because physics is 3-D.  The mathematics underlying the lattice,
however, is naturally $n$-dimensional, and the series' downstream
work consumes it at that level of generality:

\begin{itemize}
\item Paper~II's gauge-derivation programme uses lattice symmetries
      in a way that does not depend on $n = 3$;
\item the planned Hilbert-Sixth axiomatisation capstone needs
      $n$-dimensional theorems for its universality claim;
\item future kinetic-theory work on the lattice -- the Boltzmann /
      Navier-Stokes programme that Hilbert's 1900 problem singled
      out -- needs discrete-to-continuum limit machinery whose
      natural statement is dimension-agnostic.
\end{itemize}

\noindent
The mathematics is currently scattered across Paper~I's sections,
\texttt{notes/} files in upstream repositories, and the
\texttt{dcl-zoo} catalogue of the 71-dimensional per-site automorphism
algebra.  A single home prevents drift between papers and gives
downstream papers a single citable source.

\subsection{Contribution}
\label{subsec:contribution}

This paper develops the foundations of the A=1 DCL in arbitrary
dimension $n$, organised as follows.

\begin{itemize}
\item \textbf{Combinatorial foundations} (Sections~\ref{sec:lattice_n_dim}
      and~\ref{sec:automorphism_algebra}): the lattice itself in
      dimension $n$, and its automorphism algebra as a function of
      $n$ (recovering the 3-D 71-dimensional algebra at $n = 3$).
\item \textbf{Analytic foundations} (Sections~\ref{sec:discrete_diff_geometry},
      \ref{sec:function_spaces}, and~\ref{sec:discrete_measure_theory}):
      discrete differential geometry, function spaces ($\ell^p$ /
      Sobolev-analogues with $n$-dependent scaling), and discrete
      measure theory.
\item \textbf{The Hilbert-Sixth on-ramp} (Section~\ref{sec:continuum_limits}):
      discrete-to-continuum limit operators -- the most load-bearing
      content for the capstone, where the most mathematical care is
      needed.
\item \textbf{Motivating examples} (Section~\ref{sec:motivating_examples}):
      Boltzmann-shaped collision operators on the lattice, and the
      lattice-Boltzmann -> NS / Euler fluid-limit ansatz.  These
      sections are framing only at v0.1; the full proofs are out of
      scope for this first paper.
\end{itemize}

\subsection{Scope}
\label{subsec:scope}

Out of scope for this first paper:

\begin{itemize}
\item \emph{Physics applications.}  The lattice's role as the
      substrate of fundamental physics is the subject of Papers~I
      and~II and of the discrete-probability paper, not of this
      paper.
\item \emph{Sympy operationalisation.}  Machine-checked versions of
      this paper's theorems will live in the planned
      \texttt{dcl-formalism} package, not here.  This repository is
      paper-only.
\item \emph{Boltzmann / NS proof.}  The lattice kinetic theory and
      the fluid limit appear here as motivating examples
      (Section~\ref{sec:motivating_examples}).  Full proofs of the
      $H$-theorem on the lattice and of the fluid-limit convergence
      are deferred to subsequent papers.
\item \emph{Numerical experiments.}  No code accompanies this
      paper; numerical confirmation of any of its theorems happens
      either in \texttt{dcl-core}-based experiments in sibling repos
      or in the \texttt{dcl-formalism} package's automated
      sympy-verifications.
\end{itemize}

The audit table (Table~\ref{tab:audit}) is the canonical answer to
"has this claim been proved yet?".

%% The audit table is included here so the reader meets it before any
%% claim that depends on a status entry.
% ============================================================
% audit_table.tex
%
% A "system audit" table mapping the paper's claims to the evidence
% backing each one.  This is *the* canonical source of truth for which
% claims are PASS / PART / STUB / FAIL --- prose elsewhere in the paper
% should not contradict it.
%
% Status semantics (also stated in the front-matter of main.tex):
%   PASS  -- derivation complete or experiment confirms the claim to
%            stated precision.
%   PART  -- mechanism demonstrated, full quantitative match pending;
%            specific gaps documented in the evidence column.
%   STUB  -- analytical result with no full proof yet, or planned
%            proof not yet written.
%   FAIL  -- tested and disconfirmed.
%
% Use \texttt{} for code/repo names and file paths inside cells.
% Long cells are fine -- the column widths are tuned for prose.
% ============================================================

\label{tab:audit}

{\scriptsize
\setlength{\tabcolsep}{4pt}
\renewcommand{\arraystretch}{0.95}
\begin{longtable}{@{}p{2.8cm}p{3.8cm}p{3.0cm}p{4.4cm}p{1.0cm}@{}}
\caption{\small System audit: paper claims mapped to supporting
derivations and external references.  Status semantics in the
front-matter (\texttt{PASS} / \texttt{PART} / \texttt{STUB} /
\texttt{FAIL}).  Rows tagged ``inherited'' are upstream-of-record
in Paper~I and act as guardrails the diamond progression must
respect at $d = 3$.  Rows tagged ``motivating example'' describe
scope deferred to subsequent papers.}
\label{tab:audit_caption} \\
\hline
\textbf{Claim} & \textbf{Mechanism} & \textbf{Continuum / formal limit} & \textbf{Evidence} & \textbf{Status} \\
\hline
\endfirsthead

\multicolumn{5}{l}{\textit{(Table~\ref{tab:audit}, continued from previous page)}} \\
\hline
\textbf{Claim} & \textbf{Mechanism} & \textbf{Continuum / formal limit} & \textbf{Evidence} & \textbf{Status} \\
\hline
\endhead

\hline
\multicolumn{5}{r}{\textit{(continued on next page)}} \\
\endfoot

\hline
\endlastfoot

% ── Inherited guardrails from Paper I ──────────────────────────────
Bipartite octahedral lattice at $d = 3$ \emph{(guardrail, Paper~I)}
    & Joint-tick rule on the bipartite octahedral lattice
    & Carries the $(3+1)$-D Lorentz structure
    & Paper~I (\texttt{dcl}), doi:10.5281/zenodo.20078529
    & \texttt{PASS} \\

71-dim per-site automorphism algebra at $d=3$ \emph{(guardrail, \texttt{dcl-zoo})}
    & Per-site automorphisms enumerated; $71$-dim Lie algebra
    & N/A (catalogue)
    & \texttt{dcl-zoo}: generator-zoo repository
    & \texttt{PASS} \\

% ── Diamond progression construction ───────────────────────────────
Diamond progression: d-of-(d+1) simplex vertex construction
    & Combinatorial definition (Definition~\ref{def:diamond_progression}); bipartite by parity
    & Recovers Paper~I's lattice at $d = 3$ (Proposition~\ref{prop:specialisation_paper_i})
    & \S\ref{sec:lattice_n_dim} (paper-text proof pending)
    & \texttt{STUB} \\

% ── Three result-density theorems ──────────────────────────────────
Theorem~1: $\text{coord}(d) = 2d$
    & Count d-of-(d+1) simplex offsets + their negations at each site
    & Distinguishes from chemical-diamond constructions (coord bounded)
    & \S\ref{subsec:theorem_1}, Theorem~\ref{thm:coord_2d} (proof pending)
    & \texttt{STUB} \\

Theorem~2: max Euclidean hop $= \sqrt{d}$
    & Standard simplex-vertex projection into $\mathbb{R}^d$
    & Sets the rescaling parameter for $\epsilon \Lambda_d$ in the continuum limit
    & \S\ref{subsec:theorem_2}, Theorem~\ref{thm:max_hop} (proof pending)
    & \texttt{STUB} \\

Theorem~3: unit-cell volume $= (d+1)^{(d-1)/2}$
    & Simplex-volume formula applied to the offset matrix
    & Enters function-space and measure-theory scaling laws
    & \S\ref{subsec:theorem_3}, Theorem~\ref{thm:unit_cell_volume} (proof pending)
    & \texttt{STUB} \\

% ── Analytic foundations ───────────────────────────────────────────
Automorphism algebra in dimension $d$
    & Generators and commutation relations as functions of $d$; rank closed-form
    & Recovers the 3-D 71-dim per-site algebra at $d=3$
    & \S\ref{sec:automorphism_algebra}; cross-reference \texttt{dcl-zoo}
    & \texttt{STUB} \\

Discrete differential geometry on the diamond progression
    & Discrete exterior calculus: forms, $\mathrm{d}$, codifferential, Stokes; connections / curvature
    & Recovers smooth exterior calculus in the discrete-to-continuum limit (\S\ref{sec:continuum_limits})
    & \S\ref{sec:discrete_diff_geometry} (proof pending)
    & \texttt{STUB} \\

Function spaces on the diamond progression
    & Discrete $L^p$ / $\ell^p$; Sobolev-analogue $\ell^{k,p}$ with $d$-dependent scaling; embedding theorems
    & Standard $L^p$ / $W^{k,p}$ in the limit
    & \S\ref{sec:function_spaces} (proof pending)
    & \texttt{STUB} \\

Discrete measure theory on the diamond progression
    & Counting measure, Haar-analogue invariant under the automorphism algebra, change-of-variables
    & Recovers Lebesgue / Haar measure in the limit
    & \S\ref{sec:discrete_measure_theory} (proof pending)
    & \texttt{STUB} \\

Discrete-to-continuum limit operators \emph{(the Hilbert-Sixth on-ramp)}
    & Lattice spacing $\epsilon$, $\epsilon \to 0$; convergence modes (weak / strong / distributional)
    & Continuum operators in $\mathbb{R}^d$
    & \S\ref{sec:continuum_limits} (proof pending; load-bearing for the planned Hilbert-Sixth capstone)
    & \texttt{STUB} \\

% ── Headline corollary: dimensional selection ──────────────────────
Dimensional selection: $d_{\max} = 1/\delta p_{\min} - 1$
    & Per-site unity constraint at $d+1$ probability slots, each $\geq \delta p_{\min}$
    & Forces $d \leq 1/\delta p_{\min} - 1$
    & \S\ref{sec:dimensional_selection}, Theorem~\ref{thm:dp_min_dmax} (proof pending; depends on Paper~I joint-tick-rule reading)
    & \texttt{STUB} \\

$\delta p_{\min} = 1/4 \Rightarrow d_{\max} = 3$ corollary
    & Direct substitution in Theorem~\ref{thm:dp_min_dmax}
    & N/A (corollary)
    & \S\ref{subsec:dp_min_dmax_3}, Corollary~\ref{cor:dmax_3} (proof immediate once Theorem~\ref{thm:dp_min_dmax} lands)
    & \texttt{STUB} \\

% ── Motivating examples (framing only at v0.2) ─────────────────────
Boltzmann-shaped operators on the diamond progression \emph{(motivating example)}
    & Discrete collision operator; $H$-functional; equilibrium measures
    & Continuum Boltzmann equation in the limit
    & \S\ref{subsec:boltzmann} (framing only; full $H$-theorem proof deferred to a subsequent paper)
    & \texttt{STUB} \\

Fluid (NS / Euler) limit ansatz \emph{(motivating example)}
    & Lattice-Boltzmann $\to$ NS / Euler scaling
    & Navier-Stokes / Euler equations in $\mathbb{R}^d$
    & \S\ref{subsec:fluid_limit} (framing only; full convergence proof deferred)
    & \texttt{STUB} \\

\hline
\end{longtable}
}

